% chapters/01/evaluasi.tex
\bagian{Asesmen dan Instrumen Penilaian}

\begin{asesmenbox}
\textbf{Kedudukan asesmen bab ini.} Bab I tidak memiliki Sub-CPMK tersendiri
karena berperan sebagai orientasi. Asesmen pada bab ini karena itu bersifat
formatif dan tidak diberi bobot pada nilai akhir. Hasilnya dipakai dosen
untuk mengetahui kesiapan kelas dan menyesuaikan pendampingan pada Bab~III dan
seterusnya. Kerjakan kedua instrumen berikut pada minggu pertama.
\end{asesmenbox}

\subsection*{Instrumen 1: Kuis Orientasi}

\begin{enumerate}
  \item Apa perbedaan mendasar antara program permainan dan program
        penghitung nilai rata-rata yang biasa Anda tulis?
  \item Sebutkan dua capaian pembelajaran lulusan yang dibebankan pada mata
        kuliah ini dan sebutkan kode masing-masing.
  \item Berapa bobot kredit mata kuliah ini, dan berapa total beban belajar
        per minggu menurut ketentuan beban belajar?
  \item Pada minggu ke berapa mini proyek AI dikumpulkan, dan CPMK mana yang
        diukur olehnya?
  \item Sebutkan tiga syarat kelulusan mata kuliah ini.
  \item Mengapa nomor rujukan pustaka pada buku ini dibuat sama dengan nomor
        rujukan pada RPS?
\end{enumerate}

\begin{table}[htbp]
\centering
\caption{Kunci kuis orientasi}
\label{tab:kunci-kuis-orientasi}
{\setlength{\tabcolsep}{4pt}\footnotesize%
\begin{tabularx}{\textwidth}{@{}C{0.9cm} L{4.3cm} Y@{}}
\toprule
\headrow \textbf{No} & \textbf{Pokok jawaban} & \textbf{Keterangan} \\
\midrule
1 & Perulangan tanpa akhir lawan sekali jalan &
Program permainan menjalankan siklus masukan--pembaruan--penayangan secara
terus-menerus, sedangkan program penghitung berjalan sekali lalu berhenti. \\
\addlinespace
2 & \CPL{03} dan \CPL{04} &
\CPL{03} tentang perancangan dan analisis algoritma serta penerapannya pada
bahasa pemrograman; \CPL{04} tentang perancangan, implementasi, dan evaluasi
solusi berbasis \istilah{computing}. \\
\addlinespace
3 & 3 SKS; 510 menit per minggu &
Terdiri atas 150 menit tatap muka, 180 menit belajar terstruktur, dan 180
menit belajar mandiri. \\
\addlinespace
4 & Minggu 9; \CPMK{2} &
Dinilai melalui kode, laporan \istilah{playtest}, dan presentasi. \\
\addlinespace
5 & Nilai, kehadiran, keluaran wajib &
Nilai akhir minimal 60 (C); kehadiran minimal 75\%; menyerahkan mini proyek AI
dan proyek akhir. \\
\addlinespace
6 & Penelusuran antara dua dokumen &
Satu nomor rujukan menunjuk pustaka yang sama pada buku ajar dan RPS,
sehingga rujukan yang disebut di kelas dapat dicari pada kedua dokumen. \\
\bottomrule
\end{tabularx}}
\end{table}

\subsection*{Instrumen 2: Rubrik Kesiapan Repositori}

Instrumen ini menilai kesiapan tempat kerja yang akan Anda pakai sepanjang
semester. Nilainya tidak masuk ke nilai akhir, tetapi repositori yang rapi
sejak awal akan memudahkan penilaian portofolio pada minggu-minggu berikutnya.

\begin{table}[htbp]
\centering
\caption{Rubrik kesiapan repositori tugas praktikum}
\label{tab:rubrik-repositori}
{\setlength{\tabcolsep}{4pt}\footnotesize%
\begin{tabularx}{\textwidth}{@{}L{3.9cm} C{1.5cm} Y@{}}
\toprule
\headrow \textbf{Kriteria} & \textbf{Bobot} & \textbf{Deskripsi capaian} \\
\midrule
Struktur folder & 40\% &
Terdapat pemisahan yang jelas antara kode (\code{code/}), catatan
(\code{catatan/}), dan aset (\code{aset/}); penamaan berkas konsisten dan
bermakna. \\
\addlinespace
Berkas keterangan & 30\% &
Berkas \code{README.md} memuat identitas mata kuliah, nama mahasiswa, dan
deskripsi singkat rencana pekerjaan satu semester. \\
\addlinespace
Riwayat penyimpanan & 30\% &
Terdapat sekurang-kurangnya satu \istilah{commit} awal dengan pesan yang
menerangkan isinya, bukan pesan kosong atau singkatan yang tidak jelas. \\
\bottomrule
\end{tabularx}}
\end{table}

\begin{catatabox}
Apabila ketiga kriteria di atas belum terpenuhi, perbaikilah sebelum
praktikum pertama pada Bab~III. Perbaikan pada awal semester jauh lebih murah
daripada merapikan seluruh riwayat pekerjaan pada akhir semester.
\end{catatabox}
